\section{Seed-Fragility Analysis (Phase 3 Meta-Finding)}
\label{sec:fragility}

\subsection{Motivation}

Phase 1 analysis reported an ``encoder-dependent reversal'' of the EW vs NL ranking on NSL-KDD: EW $>$ NL for TargetEncoder, but NL $>$ EW for hybrid and ordinal encoders. This finding was based on 3 seeds per encoder. A natural question is whether the directional claim (EW vs NL) is robust across seeds, or whether it flips depending on the random initialization.

We investigate this by computing the sign of $\text{AUROC}_{\text{EW}} - \text{AUROC}_{\text{NL}}$ for each of 20 seeds, per dataset and encoder. The \emph{fragility rate} measures how often the direction flips:
\begin{equation}
\text{fragility rate} = \frac{\min(\text{wins}_{\text{EW}}, \text{wins}_{\text{NL}})}{n_{\text{seeds}}}
\label{eq:fragility}
\end{equation}
where $\text{wins}_{\text{EW}}$ is the number of seeds where EW $>$ NL and $\text{wins}_{\text{NL}}$ is the number where NL $>$ EW. A fragility rate of 0.0 indicates perfect agreement across all seeds; a rate of 0.5 indicates complete disagreement.

\subsection{Fragility metric}

For each (dataset, encoder) pair, we report:

\begin{itemize}
\item \textbf{EW $>$ NL count}: Number of seeds where $\text{AUROC}_{\text{EW}} > \text{AUROC}_{\text{NL}}$
\item \textbf{NL $>$ EW count}: Number of seeds where $\text{AUROC}_{\text{NL}} > \text{AUROC}_{\text{EW}}$
\item \textbf{Mean delta}: Mean of $\text{AUROC}_{\text{EW}} - \text{AUROC}_{\text{NL}}$ across seeds
\item \textbf{Std delta}: Standard deviation of the delta
\item \textbf{Fragility rate}: Proportion of seeds with the minority direction
\end{itemize}

\subsection{Results}

Table~\ref{tab:fragility} presents the seed-fragility analysis for Phase 1's target encoder on NSL-KDD, comparing the Phase 1 claim (3 seeds) with the Phase 3 replication (20 seeds).

\begin{table}[h]
\centering
\caption{Seed-fragility analysis: EW vs NL on NSL-KDD (target encoder).}
\label{tab:fragility}
\begin{tabular}{lcc}
\toprule
Metric & Phase 1 (3 seeds) & Phase 3 (20 seeds) \\
\midrule
EW $>$ NL count & 3/3 (100\%) & 4/20 (20\%) \\
NL $>$ EW count & 0/3 (0\%) & 16/20 (80\%) \\
Mean delta & +0.0390 & $-$0.024 \\
Std delta & 0.002 & 0.041 \\
Fragility rate & 0.00 & 0.20 \\
\bottomrule
\end{tabular}
\end{table}

Phase 1's claim that EW $>$ NL is fragile: the direction flips in 16 of 20 seeds. The mean delta reverses sign (from $+0.039$ in Phase 1 to $-0.024$ in Phase 3), and the standard deviation of the delta (0.041) is larger than the mean absolute effect (0.024), confirming that seed variation dominates the directional signal.

\subsection{Interpretation}

A fragility rate above 0.25 indicates that the directional claim is unreliable---the outcome depends more on which seed is used than on the underlying method comparison. For NSL-KDD with target encoder, the 20\% fragility rate (4/20 vs 16/20) suggests that while NL $>$ EW is the majority pattern, the signal is not strong enough to support confident directional claims from a single seed.

This finding is a \emph{meta-diagnostic}: it reveals that LOO's own recommendations (which method to prefer) are seed-sensitive. The practical implication is that any LOO-based conclusion drawn from a single seed should be treated as provisional until replicated across multiple seeds.

\subsection{Relationship to sign-conflict diagnostic}

The seed-fragility analysis complements the sign-conflict diagnostic from Phase 2:

\begin{itemize}
\item \textbf{Sign-conflict} (Phase 2) is a \emph{per-detector} metric: it measures how often LOO's recommendation for a single detector conflicts with held-out test effect. The 48\% rate is universal and independent of correlation.
\item \textbf{Seed-fragility} (Phase 3) is a \emph{per-direction} metric: it measures how often the overall method comparison (EW vs NL) flips across seeds. The fragility rate varies by dataset and encoder.
\end{itemize}

Both metrics support the paper's core claim: LOO-based ensemble selection is unreliable. Sign-conflict reveals that individual detector recommendations are noisy; seed-fragility reveals that aggregate method rankings are also unstable. Together, they demonstrate that LOO's diagnostic signal is unreliable at multiple levels of analysis.
